\documentclass{article}
\usepackage{../math_template}

\author{Jonathan Kasongo}
\title{Imperial-Cambridge Math Contest 2018 Rd 1 --- P2/6}

\begin{document}
\maketitle

\begin{problem}{Imperial-Cambridge Math Contest 2018 Rd 1 --- P2/6}
(i) Show that if $(k+1)$ integers are chosen from $\{1, 2, \ldots 2k+1\}$,
then among the chosen integers there are always two that are coprime.
\vspace{0.2cm}

(ii) Let $A = \{1,2,\ldots, n\}$. Prove that if $n > 11$ then there is a
bijective map $f : A \rightarrow A$ with the property that for every
$a \in A$, exactly one of $f(f(f(f(a)))) = a$ and
$f(f(f(f(f(a))))) = a$ holds.
\end{problem}

\begin{solution}{Solution}
(i) Let $1 \leq a \leq 2k$ be an integer. If both $a$ and $(a+1)$ are chosen
then we are done since $a$ and $(a+1)$ are coprime. \vspace{0.2cm}

Thus suppose we don't select two consecutive integers. There are exactly
$k$ even and $(k+1)$ odd numbers among $\{1,2,\ldots, 2k+1\}$. If we chose
the $k$ even numbers, then we would have to pick one odd number which would
create a pair of consecutive numbers. Thus the only way to pick $(k+1)$
non-consecutive numbers is to choose all of the odd numbers. However now
since 1 was chosen and every number is coprime with 1, there does infact
exist two coprime numbers among the chosen integers. \\

(ii) We examine two cases. \vspace{0.2cm}

\textit{Case 1: $n$ is even.} Then define $f$ like so,
$$
f(a) =
\begin{cases}
a+1 \quad \text{if } a \text{ is odd},\\
a-1 \quad \text{if } a \text{ is even}.
\end{cases}
$$
For odd $a$ we have $f(f(a)) = a$ and so $f(f(f(f(a)))) = a$,
whilst $f(f(f(f(f(a))))) = f(a) = a+1$. Similarly for even $a$ we have
$f(f(a)) = a$ and so $f(f(f(f(a)))) = a$ whilst
$f(f(f(f(f(a))))) = f(a) = a-1$, as desired.\vspace{0.2cm}

\textit{Case 2: $n$ is odd.} Then define $f$ like so,
$$
f(n-4) = n-3, \ \  f(n-3) = n-2, \ \ f(n-2) = n-1, \ \ f(n-1) = n, \ \
f(n) = n-4
$$
and for $1 \leq a < n-4$,
$$
f(a) =
\begin{cases}
a+1 \quad \text{if } a \text{ is odd},\\
a-1 \quad \text{if } a \text{ is even}.
\end{cases}
$$
By the same arguement as in case 1, clearly the desired property is true
for $1 \leq a < n-4$. But now for each $n-4 \leq a \leq n$ it can easily
be verified that only $f(f(f(f(f(a))))) = a$ holds. $\blacksquare$
\end{solution}

\end{document}
